% Part: proof-theory
% Chapter: normalization
% Section: introduction

\documentclass[../../../include/open-logic-section]{subfiles}

\begin{document}

\olfileid{pt}{nor}{int}

\olsection{مقدمه}

اگر شرطی~$!A\lif !B$ را !!{prove} کرده باشید، شاید بتوانید با
\Elim{\lif} از آن در !!a{proof} برای~$!B$ استفاده کنید. البته برای این
کار به !!a{proof}~$\delta_1$ برای~$!A$ نیز نیاز دارید. !!{proof} شما
برای~$!A\lif !B$ به احتمال زیاد به~\Intro{\lif} پایان می‌یابد؛ قاعده‌ای
که !!a{proof}~$\delta_2$ برای~$!B$ از فرض باز~$!A$ را به برهانی برای
$!A\lif !B$ تبدیل می‌کند. اگر چنین باشد، !!{proof} نهایی شما این شکل
را دارد:
\[
\AxiomC{$\Discharge{!A}{x}$}
\RightLabel{$\delta_2$}
\DeduceC{$!B$}
\DischargeRule{\Intro{\lif}}{x}
\UnaryInfC{$!A \lif !B$}
\AxiomC{}
\RightLabel{$\delta_1$}
\DeduceC{$!A$}
\RightLabel{\Elim{\lif}}
\BinaryInfC{$!B$}
\DisplayProof
\]
راهبرد شما از این جهت کارآمد است که چیزی را که از پیش داشتید، یعنی
!!a{proof}ی برای~$!A\lif !B$، دوباره به کار بردید؛ اما خود !!{proof}
کارآمد نیست. معرفی $!A\lif !B$ فقط برای حذف بی‌درنگ آن، راهی انحرافی
به نظر می‌رسد. باید راه مستقیم‌تری برای !!{prove} کردن~$!B$ وجود داشته
باشد.

در واقع همیشه چنین راهی هست. «انحراف»هایی مانند بالا---!!a{introduction}
که بی‌درنگ !!a{elimination} در پی آن می‌آید---همیشه از
!!{proof}های استنتاج طبیعی حذف‌شدنی‌اند. این فرایند \emph{نرمال‌سازی}
نام دارد و حاصل آن یک !!{proof} \emph{نرمال}، یا !!a{proof}
\emph{در صورت نرمال}، است.

اثبات این نتیجه، مانند قضیهٔ حذف برش در حساب دنباله‌ای، تا اندازه‌ای
درگیر است و بررسی حالت‌های فراوانی می‌طلبد. اثبات آن برای \Log{N1c}
کلاسیک دشوارتر از \Log{N1i} شهودگرایانه است. همان‌گونه که حذف برش
نشان می‌دهد با قاعدهٔ \Cut{} در حساب دنباله‌ای چیز بیشتری را نمی‌توان
!!{prove} کرد، نرمال‌سازی نشان می‌دهد استفاده از انحراف‌ها چیزی بیش از
پرهیز از آن‌ها اثبات نمی‌کند. شباهت‌های دیگری نیز هست: !!{proof} نرمال
ممکن است بسیار بلندتر از کوتاه‌ترین !!{proof} دارای انحراف برای همان
نتیجه باشد، درست همان‌گونه که !!{proof}های حساب دنباله‌ای با برش
می‌توانند بسیار کوتاه‌تر از کوتاه‌ترین !!{proof} بدون برش باشند.
خاصیت زیر-!!{formula}، یعنی اینکه همواره می‌توان نتیجه‌ای را !!{prove}
کرد و در این کار تنها زیر-!!{formula}های آن نتیجه و فرض‌های بازِ
!!{proof} را به کار برد،
برای استنتاج طبیعی از نرمال‌سازی نتیجه می‌شود، همان‌گونه که برای حساب
دنباله‌ای از حذف برش نتیجه می‌شود.

\end{document}
